\documentclass[12pt]{article}
\usepackage[a4paper,margin=1.6cm]{geometry}
\usepackage{amsmath,amssymb}
\usepackage{tasks}
\usepackage{multicol}
\usepackage{array}
\usepackage{fancyhdr}
\usepackage{enumitem}
\usepackage{setspace}
\setlength{\parindent}{0pt}
\setlength{\parskip}{6pt}
\renewcommand{\familydefault}{\sfdefault}

\newcommand{\blank}{\rule{3.3cm}{0.4pt}}
\newcommand{\blanklong}{\rule{7.2cm}{0.4pt}}
\newcommand{\tickbox}{\(\square\)}
\newcommand{\overbar}[1]{\overline{#1}}

\pagestyle{fancy}
\fancyhf{}
\lhead{\textbf{LCA Maths} \;—\; Fractions, Decimals \& Percentages}
\rhead{\textbf{Name:} \blank\;\;\textbf{Date:} \blank}
\cfoot{\thepage}

\begin{document}
\Large

{\LARGE \textbf{Flexibly convert between fractions, decimals, and percentages}}\\[2pt]
{\Large Evidence of learning worksheet (after 20 minutes using the app)}\\

\textbf{Rules for this worksheet (use your calculator):}
\begin{itemize}[itemsep=2pt,topsep=2pt]
\item If a decimal \textbf{terminates}, write it exactly (e.g.\ $0.125$).
\item If a decimal \textbf{repeats}, write it using \textbf{$\ldots$} (e.g.\ $0.111\ldots$) \textit{or} with a bar if you know it (e.g.\ $0.\overbar{1}$).
\item For percentages that are not exact, round to the \textbf{nearest whole \%} unless told otherwise.
\end{itemize}

\hrule
\vspace{6pt}

% =============== SECTION 1 ===============
{\Large \textbf{1) Warm-up: fill the missing form}}\\
Complete each line. (Write \textbf{all three forms}.)\\[-2pt]
\begin{tasks}(2)
\task \( \dfrac{3}{4} = \blank \) as a decimal,\quad \(=\blank\%\)
\task \( 0.6 = \blank \) as a fraction,\quad \(=\blank\%\)
\task \( 25\% = \blank \) as a fraction,\quad \(=\blank \) as a decimal
\task \( 0.125 = \blank \) as a fraction,\quad \(=\blank\%\)
\task \( \dfrac{7}{20} = \blank \) as a decimal,\quad \(=\blank\%\)
\task \( 0.45 = \blank \) as a fraction,\quad \(=\blank\%\)
\end{tasks}

\vspace{4pt}
% =============== SECTION 2 ===============
{\Large \textbf{2) Terminating or repeating?}}\\
For each fraction, \textbf{tick} whether the decimal \textbf{terminates} or \textbf{repeats}.\\[-2pt]
\renewcommand{\arraystretch}{1.4}
\begin{tabular}{|>{\centering\arraybackslash}m{2.6cm}|>{\centering\arraybackslash}m{3.3cm}|>{\centering\arraybackslash}m{3.3cm}|>{\arraybackslash}m{6.2cm}|}
\hline
\textbf{Fraction} & \textbf{Terminates} & \textbf{Repeats} & \textbf{Reason (short)}\\
\hline
\(\dfrac{3}{8}\)  & \tickbox & \tickbox & \blanklong\\ \hline
\(\dfrac{5}{12}\) & \tickbox & \tickbox & \blanklong\\ \hline
\(\dfrac{7}{25}\) & \tickbox & \tickbox & \blanklong\\ \hline
\(\dfrac{11}{40}\)& \tickbox & \tickbox & \blanklong\\ \hline
\(\dfrac{2}{9}\)  & \tickbox & \tickbox & \blanklong\\ \hline
\(\dfrac{7}{11}\) & \tickbox & \tickbox & \blanklong\\ \hline
\end{tabular}

\vspace{6pt}
% =============== SECTION 3 ===============
{\Large \textbf{3) Ninths (denominator 9) — do the full set}}\\
Complete the table. Write the decimal using \(\ldots\) and round the percent to the nearest whole \%.\\[-4pt]
\renewcommand{\arraystretch}{1.45}
\begin{tabular}{|>{\centering\arraybackslash}m{2.2cm}|>{\centering\arraybackslash}m{5.2cm}|>{\centering\arraybackslash}m{4.0cm}|}
\hline
\textbf{Fraction} & \textbf{Decimal} & \textbf{Percent (nearest \%)}\\
\hline
\(\dfrac{1}{9}\) & \blank & \blank\\ \hline
\(\dfrac{2}{9}\) & \blank & \blank\\ \hline
\(\dfrac{3}{9}\) & \blank & \blank\\ \hline
\(\dfrac{4}{9}\) & \blank & \blank\\ \hline
\(\dfrac{5}{9}\) & \blank & \blank\\ \hline
\(\dfrac{6}{9}\) & \blank & \blank\\ \hline
\(\dfrac{7}{9}\) & \blank & \blank\\ \hline
\(\dfrac{8}{9}\) & \blank & \blank\\ \hline
\end{tabular}

\vspace{6pt}
% =============== SECTION 4 ===============
{\Large \textbf{4) Elevenths (denominator 11) — do the full set}}\\
Complete the table. Write the repeating decimal using \(\ldots\). Round the percent to the nearest whole \%.\\[-4pt]
\renewcommand{\arraystretch}{1.45}
\begin{tabular}{|>{\centering\arraybackslash}m{2.2cm}|>{\centering\arraybackslash}m{5.2cm}|>{\centering\arraybackslash}m{4.0cm}|}
\hline
\textbf{Fraction} & \textbf{Decimal} & \textbf{Percent (nearest \%)}\\
\hline
\(\dfrac{1}{11}\) & \blank & \blank\\ \hline
\(\dfrac{2}{11}\) & \blank & \blank\\ \hline
\(\dfrac{3}{11}\) & \blank & \blank\\ \hline
\(\dfrac{4}{11}\) & \blank & \blank\\ \hline
\(\dfrac{5}{11}\) & \blank & \blank\\ \hline
\(\dfrac{6}{11}\) & \blank & \blank\\ \hline
\(\dfrac{7}{11}\) & \blank & \blank\\ \hline
\(\dfrac{8}{11}\) & \blank & \blank\\ \hline
\(\dfrac{9}{11}\) & \blank & \blank\\ \hline
\(\dfrac{10}{11}\) & \blank & \blank\\ \hline
\end{tabular}

\vspace{6pt}
% =============== SECTION 5 ===============
{\Large \textbf{5) Fix the common mix-up (ninths vs elevenths)}}\\
Answer each question clearly.

\begin{enumerate}[label=\textbf{(\alph*)},itemsep=4pt]
\item Which is closer to \(78\%\): \(\dfrac{7}{9}\) or \(\dfrac{7}{11}\)? Show your working.\\
\blanklong\\ \blanklong

\item Fill in the missing values (nearest whole \%):\\[-4pt]
\begin{tasks}(2)
\task \( \dfrac{7}{9} \approx \blank\% \)
\task \( \dfrac{7}{11} \approx \blank\% \)
\end{tasks}

\item Explain (one sentence): why do ninths give \(0.111\ldots\) style decimals?\\
\blanklong
\end{enumerate}

\vspace{6pt}
% =============== SECTION 6 ===============
{\Large \textbf{6) Mixed conversions (not all are nice)}}\\
Complete each one. (Simplify fractions.)

\begin{tasks}(2)
\task \( 0.3 = \blank \) as a fraction,\quad \(=\blank\%\)
\task \( 0.08 = \blank \) as a fraction,\quad \(=\blank\%\)
\task \( 62.5\% = \blank \) as a fraction,\quad \(=\blank \) as a decimal
\task \( 7.5\% = \blank \) as a fraction,\quad \(=\blank \) as a decimal
\task \( 0.09\ldots = \blank \) as a fraction
\task \( 0.18\ldots = \blank \) as a fraction
\end{tasks}

\vspace{6pt}
% =============== SECTION 7 ===============
{\Large \textbf{7) Challenge: choose the best match}}\\
Each percentage matches \textbf{one} fraction. Write the correct letter beside each.\\[-6pt]

\begin{multicols}{2}
\textbf{Fractions:}\\[-6pt]
\begin{enumerate}[label=\textbf{\Alph*.},itemsep=2pt]
\item \(\dfrac{4}{9}\)
\item \(\dfrac{7}{9}\)
\item \(\dfrac{7}{11}\)
\item \(\dfrac{5}{8}\)
\item \(\dfrac{3}{20}\)
\item \(\dfrac{9}{40}\)
\end{enumerate}

\columnbreak

\textbf{Percents:}\\[-6pt]
\begin{enumerate}[label=\textbf{1.},itemsep=2pt]
\item \(44\%\) \;\;\(\rightarrow\)\; \blank
\item \(78\%\) \;\;\(\rightarrow\)\; \blank
\item \(64\%\) \;\;\(\rightarrow\)\; \blank
\item \(63\%\) \;\;\(\rightarrow\)\; \blank
\item \(15\%\) \;\;\(\rightarrow\)\; \blank
\item \(23\%\) \;\;\(\rightarrow\)\; \blank
\end{enumerate}
\end{multicols}

\vspace{6pt}
\hrule
\vspace{4pt}
{\Large \textbf{Reflection (2 lines):}}\\
After using the app and doing this worksheet, what is one pattern you noticed for ninths or elevenths?\\
\blanklong\\
\blanklong

\end{document}
